\documentclass[11pt,a4paper]{article}
\usepackage[utf8]{inputenc}
\usepackage{amsmath,amssymb,amsthm}
\usepackage{listings}
\usepackage[margin=2.5cm]{geometry}
\usepackage{xcolor}

\lstset{
  language=C++,
  basicstyle=\ttfamily\small,
  keywordstyle=\color{blue}\bfseries,
  commentstyle=\color{green!50!black},
  stringstyle=\color{red!70!black},
  numbers=left,
  numberstyle=\tiny\color{gray},
  breaklines=true,
  frame=single,
  tabsize=4,
  showstringspaces=false
}

\title{IOI 1991 -- Problem 3: Matrix Game}
\author{Solution}
\date{}

\begin{document}
\maketitle

\section{Problem Statement}

Two players play on an $m \times n$ matrix of integers. They alternate turns.
On each turn, a player selects an entire remaining row or column, receives
the sum of its entries, and removes it from the matrix. The game ends when no
rows or columns remain. Player~1 maximizes their total score; Player~2
minimizes Player~1's total (equivalently, maximizes their own).

Determine Player~1's optimal score under perfect play.

\section{Solution}

\subsection{Minimax with Bitmask Memoization}

This is a two-player zero-sum game solved by minimax:

\begin{itemize}
  \item \textbf{State:} A pair of bitmasks $(\textit{rowMask},
    \textit{colMask})$ indicating which rows and columns remain, plus a
    boolean for whose turn it is.
  \item \textbf{Transitions:} Remove one remaining row or column.
  \item \textbf{Value:} Player~1 maximizes accumulated score; Player~2
    minimizes it.
\end{itemize}

The whose-turn information is determined by the parity of the number of
removed rows/columns, so it need not be stored separately (though we include
it for clarity).

\subsection{State Space}

With bitmasks for rows ($2^m$ states) and columns ($2^n$ states), plus the
turn bit, the total number of states is $2 \times 2^m \times 2^n$. For
$m, n \le 10$, this is $2^{21} \approx 2 \times 10^6$, which is feasible.

\subsection{Bug Fix: Hash Key Encoding}

The original code used the encoding
$\textit{key} = (\textit{rowMask} \ll 20) \mid (\textit{colMask} \ll 8)
\mid \textit{isMax}$, which causes collisions when $n > 8$ (since
$\textit{colMask}$ can exceed 8 bits). The corrected version below uses a
collision-free encoding.

\section{Complexity Analysis}

\begin{itemize}
  \item \textbf{Time:} $O(2^{m+n} \cdot (m+n) \cdot \max(m,n))$. Each of
    the $O(2^{m+n})$ states has $O(m+n)$ transitions, each requiring
    $O(\max(m,n))$ to compute the row/column sum.
  \item \textbf{Space:} $O(2^{m+n})$ for memoization.
\end{itemize}

For $m, n \le 10$, this is comfortably within time limits.

\section{Implementation}

\begin{lstlisting}
#include <iostream>
#include <unordered_map>
#include <climits>
#include <algorithm>
using namespace std;

int m, n;
int A[12][12];
unordered_map<long long, int> memo;

int rowSum(int r, int colMask) {
    int s = 0;
    for (int c = 0; c < n; c++)
        if (colMask & (1 << c))
            s += A[r][c];
    return s;
}

int colSum(int c, int rowMask) {
    int s = 0;
    for (int r = 0; r < m; r++)
        if (rowMask & (1 << r))
            s += A[r][c];
    return s;
}

// Returns Player 1's optimal score from this state onward.
int solve(int rowMask, int colMask, bool isMax) {
    if (rowMask == 0 || colMask == 0) return 0;

    // Collision-free key: rowMask needs m bits, colMask needs n bits
    long long key = ((long long)isMax << 30) |
                    ((long long)rowMask << 15) |
                    (long long)colMask;
    auto it = memo.find(key);
    if (it != memo.end()) return it->second;

    int best = isMax ? INT_MIN : INT_MAX;

    // Try removing each remaining row
    for (int r = 0; r < m; r++) {
        if (!(rowMask & (1 << r))) continue;
        int s = rowSum(r, colMask);
        int future = solve(rowMask ^ (1 << r), colMask, !isMax);
        if (isMax)
            best = max(best, s + future);
        else
            best = min(best, future);
    }

    // Try removing each remaining column
    for (int c = 0; c < n; c++) {
        if (!(colMask & (1 << c))) continue;
        int s = colSum(c, rowMask);
        int future = solve(rowMask, colMask ^ (1 << c), !isMax);
        if (isMax)
            best = max(best, s + future);
        else
            best = min(best, future);
    }

    memo[key] = best;
    return best;
}

int main() {
    cin >> m >> n;
    for (int i = 0; i < m; i++)
        for (int j = 0; j < n; j++)
            cin >> A[i][j];

    int result = solve((1 << m) - 1, (1 << n) - 1, true);
    cout << result << endl;
    return 0;
}
\end{lstlisting}

\section{Notes}

\begin{itemize}
  \item The minimax approach gives the exact optimal result under perfect
    play from both players.
  \item Alpha-beta pruning can speed up the search in practice but does not
    improve worst-case complexity.
  \item The memoization key must encode \textit{rowMask}, \textit{colMask},
    and the turn bit without collisions. With $m, n \le 12$, using 15 bits
    each for the masks and 1 bit for the turn suffices within a
    \texttt{long long}.
\end{itemize}

\end{document}
